\subsection{标量的限制与扩张}

设 \(A_{0}\) 和 A 是两个交换环，且 \(\rho:A_{0}\to A\) 是一个环同态。若 E 是一个 A-代数，我们（与 II，§ 1，第 13 号相一致地）用 \(\rho_{*}(E)\) 表示 \(A_{0}\) -模，它由 E 上的加法与外合成律

\[
\lambda . x = \rho (\lambda) x \quad \text { for   all } \lambda \in A _ {0} \text { and   all } x \in E.
\]

E 中的乘法与 \(A_{0}\) -模结构（在 \(\rho_{*}(E)\) 上）一起定义了一个 \(A_{0}\) -代数结构（在 \(\rho_{*}(E)\) 上）。当 \(A_{0}\) 是 A 的子环且 \(\rho\) 为典范单射时，称代数 \(\rho_{*}(E)\) 是由 E 把标量环 A 限制到 \(A_{0}\) 而得到的。作为语言上的滥用，当 \(\rho\) 任意时有时也这样说。

设 F 为一个 \(A_{0}\) -代数。 \(A_{0}\) -代数同态 \(F \to \rho_{*}(E)\) 称为半同态（相对于 \(\rho\) ），或 F 到 A-代数 E 中的 \(\rho\) -同态；若不致混淆，也称为 \(A_{0}\) -同态。若 E, \(E'\) 是两个 A-代数，则每个 A-代数同态 \(E \to E'\) 也是 \(A_{0}\) -代数同态 \(\rho_{*}(E) \to \rho_{*}(E')\) 。

现在考虑两个交换环 A 和 B 以及一个环同态 \(\rho: \mathbf{A} \to \mathbf{B}\) 。对每个 A-模 E，由 E 把标量环 A 扩张到 B 而得到的 B-模 \(\rho^{*}(\mathbf{E}) = \mathbf{E} \otimes_{\mathbf{A}} \mathbf{B}\) 已经定义（II，§ 5，第 1 号）。若 E 还是一个 A-代数，我们要在 \(\rho^{*}(\mathbf{E})\) 上定义一个 B-代数结构。为此，注意到 \((\mathbf{E} \otimes_{\mathbf{A}} \mathbf{B}) \otimes_{\mathbf{B}} (\mathbf{E} \otimes_{\mathbf{A}} \mathbf{B})\) 典范同构于 \((\mathbf{E} \otimes_{\mathbf{A}} \mathbf{E}) \otimes_{\mathbf{A}} \mathbf{B}\) （II，§ 5，第 1 号，命题 3）。若 \(m: \mathbf{E} \otimes_{\mathbf{A}} \mathbf{E} \to \mathbf{E}\) 定义 E 上的乘法，则映射 \(m \otimes l_{\mathbf{B}}: (\mathbf{E} \otimes_{\mathbf{A}} \mathbf{E}) \otimes_{\mathbf{A}} \mathbf{B} \to \mathbf{E} \otimes_{\mathbf{A}} \mathbf{B}\) 便典范地等同于一个 B-线性映射

\[
m ^ {\prime}: \rho^ {*} (\mathbf {E}) \otimes_ {\mathbf {B}} \rho^ {*} (\mathbf {E}) \rightarrow \rho^ {*} (\mathbf {E})
\]

它定义了 \(\rho^{*}(\mathbf{E})\) 上所要的 B-代数结构。于是

\[
(x \otimes \beta) (x ^ {\prime} \otimes \beta^ {\prime}) = (x x ^ {\prime}) \otimes (\beta \beta^ {\prime})\tag{5}
\]

对 \(x, x'\) （属于 \(\mathbf{E}\) ）与 \(\beta\) 、 \(\beta'\) （属于 \(\mathbf{B}\) ）成立，这便是所要的 B-代数结构。称 B-代数 \(\rho^*(\mathbf{E})\) 为由 A-代数 \(\mathbf{E}\) 通过把标量环 \(\mathbf{A}\) 扩张到 \(\mathbf{B}\) （借助 \(\rho\) ）而得到的代数。它也记作 \(\mathrm{E}_{(\mathbf{B})}\) 或 \(\mathbf{E} \otimes_{\mathbf{A}} \mathbf{B}\) 。当 \(\mathbf{E}\) 为结合的（相应地，交换的，相应地，含单位元的）时， \(\rho^*(\mathbf{E})\) 亦然。

对每个 A-代数 E，典范映射 \(\phi_{\mathbf{E}}: x \mapsto x \otimes 1\) （由 E 到 \(\mathrm{E}_{(\mathbf{B})}\) ）是一个代数的 A-同态。此外，对每个 B-代数 F 和每个 A-同态 \(f: \mathbf{E} \to \mathbf{F}\) ，存在唯一的 B-同态 \(\bar{f}: \mathbf{E}_{(\mathbf{B})} \to \mathbf{F}\) ，使得 \(\bar{f}(x \otimes 1) = f(x)\) 对所有 \(x \in \mathbf{E}\) 成立。

第一个断言由 \(\mathrm{E}_{(\mathrm{B})}\) 中乘法的定义直接得出，该定义给出 \((x\otimes1)(x'\otimes1)=(xx')\otimes1\) （对 \(x\in E\) 与 \(x'\in E\) ）。 B-线性映射 \(\bar{f}\) ——它由 \(\mathrm{E}_{(\mathrm{B})}\) 到 F 中、满足关系 \(\bar{f}(x\otimes1)=f(x)\) （对所有 \(x\in E\) ）——的存在性与唯一性由 II，§ 5，第 1 号，命题 1 得出；于是只需验证 \(\bar{f}(yy')=\bar{f}(y)\bar{f}(y')\) 对 y 与 \(y'\) （ \(\mathrm{E}_{(\mathrm{B})}\) 中的）成立；由于形如 \(x\otimes1\) （其中 \(x\in E\) ）的元素生成 B-模 \(\mathrm{E}_{(\mathrm{B})}\) ，只需考虑 \(y=x\otimes1\) , \(y'=x'\otimes1\) （其中 \(x\in E\) , \(x'\in E\) ）的情形；由于 \(yy'=(xx')\otimes1\) ，关系 \(\bar{f}(yy')=\bar{f}(y)\bar{f}(y')\) 便由 \(f(xx')=f(x)f(x')\) 得出。

也可以说 \(f \mapsto \bar{f}\) 是一个典范双射

\[
\operatorname{Hom} _ {\mathrm{A-alg.}} (\mathrm{E}, \rho_ {*} (\mathrm{F})) \rightarrow \operatorname{Hom} _ {\mathrm{B-alg.}} (\rho^ {*} (\mathrm{E}), \mathrm{F}).\tag{6}
\]

由 \(E_{(B)}\) 与 \(\phi_{E}\) 组成的有序对因此是泛映射问题（《集合论》，IV，§ 3，第 1 号）的一个解，其中 \(\Sigma\) 是 B-代数结构这个种类，而 \(\alpha\) -映射是 E 到 B-代数中的 A-同态。

推论。设 E, E' 为两个 A-代数；对代数的每个 A-同态 \(u: E \to E'\) ， \(u \otimes l_{B}\) 是使图表交换的唯一的代数 B-同态 v: \(E \otimes_{A} B \to E' \otimes_{A} B\)

\[
\begin{array}{c} \mathrm{E} \xrightarrow {\phi_ {\mathrm{E}}} \mathrm{E} \otimes_ {\mathrm{A}} \mathrm{B} \\ u \Biggl \downarrow \\ \mathrm{E} ^ {\prime} \xrightarrow [ \phi_ {\mathrm{E} ^ {\prime}} ]{} \mathrm{E} ^ {\prime} \otimes_ {\mathrm{A}} \mathrm{B} \end{array}
\]

设 C 为第三个交换环，且 \(\sigma: B \to C\) 为一个环同态；立即可见，典范 C-同态

\[
\sigma^ {*} (\rho^ {*} (\mathbf {E})) \rightarrow (\sigma \circ \rho) ^ {*} (\mathbf {E})
\]

把 \((x\otimes 1)\otimes 1\) 映为 \(x\otimes 1\) （对所有 \(x\in \mathbf{E}\) ）者（II，§ 5，第 1 号，命题 2）是一个代数同构。

